\section*{Análisis de Riesgo}
\addcontentsline{toc}{section}{Análisis de Riesgo}

Para la creación de la compañía, es necesario que los procesos de gestión y planificación cuenten con un análisis de riesgos, ya que este permite identificar los riesgos que puedan afectar negativamente a la empresa o a los que pueda estar expuesta. Por esta razón, es fundamental establecer estrategias que permitan mitigar dichos riesgos o reducir la probabilidad de que ocurran.

El análisis de riesgos busca examinar la magnitud y naturaleza de los posibles efectos negativos derivados de la propuesta, así como la probabilidad de que estos se materialicen. Del mismo modo, permite identificar medidas eficaces para reducir dichos riesgos y evaluar alternativas a la propuesta planteada.

\subsection*{Factores limitantes}
\addcontentsline{toc}{subsection}{Factores limitantes}

Los factores limitantes son aquellos elementos o circunstancias que impiden el progreso y el logro de los objetivos establecidos en la empresa. Estos pueden ser internos, como la falta de recursos o capacidades del equipo, o externos, como la competencia intensa o cambios en el entorno económico. 

En el contexto del proyecto, estos factores representan barreras clave que deben gestionarse para garantizar la viabilidad del negocio

\begin{enumerate}
	\item \textbf{Costo de adquirir usuarios:} Como plataforma emergente en el ecosistema edtech bogotano, requeriremos una inversión significativa en marketing digital, incluyendo campañas en redes sociales. Por lo que los costos iniciales por usuario adquirido podrían ser elevados si las campañas no logran una alta efectividad.
	
	\item \textbf{Dependencia de oferta inicial:} La atractividad de la plataforma depende de una oferta inicial robusta de usuarios para fomentar interacciones comunitarias y generar datos para el machine learning. Sin un número suficiente de usuarios activos, será difícil crear el impulso viral necesario, lo que podría llevar a una adopción lenta y comprometer la sostenibilidad del modelo freemium.
	
	\item \textbf{Competencia intensa:} Competidores establecidos dominan el mercado con servicios de control financiero y banca digital. Existe el riesgo de que lancen módulos educativos similares, exigiendo una diferenciación continua y reforzar la propuesta de valor con el tiempo.
	
	\item \textbf{Brecha digital y demográfica:} Aunque Bogotá tiene alta penetración digital, persisten barreras en estratos bajos y trabajadores informales, con informalidad laboral alta, lo que podría limitar en cierto aspecto la adopción inicial y requiriendo estrategias de onboarding inclusivas.
\end{enumerate}

\subsection*{Factores claves de éxito}
\addcontentsline{toc}{subsection}{Factores claves de éxito}

Los factores clave de éxito (FCE) son las capacidades y características del producto altamente valoradas por los clientes, en las que la empresa debe tener éxito para superar a competidores y asegurar sostenibilidad a largo plazo. La identificación de estos factores dentro de una compañía es esencial para comprender cuales son los aspectos que colaboran al rendimiento superior y sostenibilidad de la empresa en el largo plazo. 

Algunos de estos factores que se pueden presentar en nuestra empresa son los siguientes:

\begin{enumerate}
	\item \textbf{Estrategia de marketing digital:} Desarrollar campañas de marketing efectivas dirigidas principalmente a nuestro público objetivo mediante redes sociales, influencers y publicidad digital, esto será esencial para atraer a personas interesadas en este proyecto. La estrategia deberá destacar la facilidad de uso, el manejo de las recomendaciones de inversión, el manejo de elementos audiovisuales y la personalización via machine learning.
	
	\item \textbf{Construcción de comunidad:} Fomentar el inicio de una comunidad activa de usuarios que puedan interactuar dentro y fuera de la plataforma, compartir sus experiencias personales y recomendarla, esto será la base para poder construir fundamentos sólidos y generar lealtad en los usuarios.
	
	\item \textbf{Escalabilidad técnica:} Asegurar que la plataforma logre ser escalable desde el punto de vista tecnológico haciendo que esta pueda crecer sin tener complicaciones desde la operación o técnica a medida que se aumente la cantidad de usuarios. Optimizando y aprovechando la arquitectura de microservicios para soportar crecimiento exponencial de usuarios sin caídas.
	
	\item \textbf{Calidad de la plataforma:} Permitir usar la plataforma con anuncios en su versión gratis, buscando asegurar buenas impresiones en nuestros usuarios de forma que decidan pagar la suscripción para remover anuncios y desbloquear mayores funcionalidades.
	
	\item \textbf{Validación institucional:} Realizar alianzas con entidades como Superintendencia Financiera y programas como Banca de las Oportunidades para credibilidad y co-marketing.
\end{enumerate}


\subsection*{Riesgos legales}
\addcontentsline{toc}{subsection}{Riesgos legales}

Los riesgos legales hablan sobre las posibles consecuencias legales que la compañía podría afrontar como efecto de la toma o desconocimiento de acciones que se interpretan como ilegales o inapropiadas. Estos riesgos se manifiestan en varias situaciones, como por ejemplo las relaciones laborales, cumplimiento de regulaciones vigentes, problemas con los derechos de propiedad intelectual, problemas fiscales, seguridad de la información y muchas otras áreas más.

Las repercusiones de estos escenarios pueden generar la amenaza de demandas civiles, imposición de multas o sanciones por parte de las entidades gubernamentales, e inclusive, en el peor de los casos, poder enfrentar cargos penales. 

Para nuestra plataforma, se identifican los siguientes riesgos:

\begin{enumerate}
	\item \textbf{Protección de datos personales:} Debido a que es un servicio de planificación financiera, la empresa va a manejar información personal financiera de usuarios que podría ser altamente sensible. Además de métodos de pago con el tema de las suscripciones. Por tanto, para mitigar estos riesgos, se debe asegurar el cumplimiento de la protección de datos, la realización de auditorías regulares, establecer los consentimientos de la plataforma con el usuario, y establecer protocolos de seguridad informática robustos.
	
	\item \textbf{Cumplimiento normativo:} Debido a ser una plataforma edtech con ciertas características de una fintech, estamos sujetos a seguir los registros dados por la Superintendencia Financiera y cumplir con las regulaciones emergentes del ecosistema fintech, para moderar estos riesgos, nos aseguraremos de implementar correctamente los registros ante las respectivas autoridades pertinentes, además de consultar a la Superintendencia y otras entidades en caso de dudas sobre la autorización, vigilancia o registro en específico de cierto tipo de actividades.
	
	\item \textbf{Derechos de propiedad:} Implicación de posibles conflictos generados por el manejo inadecuado de la protección de derechos de autor de los desarrollos tecnológicos propios (como el contenido educativo y los algoritmos de machine learning) además de la implementación de derechos de terceros. Por esto, debemos establecer políticas claras para no alterar aquellos derechos, entregar créditos a los respectivos autores y consultas sobre registro de derechos de autor ante la Superintendencia.
\end{enumerate}


\subsection*{Riesgos operacionales}
\addcontentsline{toc}{subsection}{Riesgos operacionales}

La gestión de riesgos operacionales es un proceso cíclico continuo que busca identificar, evaluar y controlar los riesgos relacionados a operaciones diarias de la empresa, lo que conlleva a su aceptación, mitigación o elusión.Esta gestión implica la supervisión constante del riesgo operativo, incluyendo la posibilidad de pérdidas derivadas de fallas en los procesos y sistemas internos, factores humanos u otros eventos externos \parencite{Microbank01}.

\begin{enumerate}
	\item \textbf{Fallas técnicas y de vulnerabilidades:} Se refiere a la necesidad de que la infraestructura de la plataforma sea lo suficientemente robusta para evitar caídas del sistema o problemas de rendimiento, especialmente en momentos de alta demanda, así como la aparición de vulnerabilidades. Para ello, se implementará una infraestructura sólida, junto con protocolos de ciberseguridad y la realización de pruebas constantes.
	
	\item \textbf{Gestión del crecimiento:} A medida que se expandan las operaciones de la compañía, un crecimiento acelerado y no planificado podría generar dificultades en la gestión de la plataforma. Tanto la atención al cliente como la calidad del servicio podrían verse afectadas. Por esta razón, se propone llevar a cabo una expansión gradual, planificada y controlada de las operaciones hacia nuevos mercados.
	
	\item \textbf{Incumplimiento de objetivos:} Se refiere a la incapacidad de la empresa para cumplir con los objetivos y plazos establecidos en el desarrollo de nuevas iniciativas o proyectos. Para mitigar este riesgo, se realizará un estudio detallado de los proyectos que ingresen a la empresa, así como un análisis adecuado de los recursos humanos involucrados en su ejecución.
\end{enumerate}


